\section{What the Church Actually Does}

Doctrine about death can be tested against a body of practice that is unusually well documented, because most of it is law. What the Church does at a deathbed is prescribed, the prescriptions have been revised twice in living memory, and the revisions were made in explicit reference to a conciliar text of 1551. It is a rare case where a reader can check whether the practice matches the teaching, and the answer is instructive: it matches the teaching better than it matches what most Catholics believe the teaching to be.

\subsection{The apostolic text, read where it stands}

``Is any man sick among you? Let him bring in the priests of the church and let them pray over him, anointing him with oil in the name of the Lord. And the prayer of faith shall save the sick man. And the Lord shall raise him up: and if he be in sins, they shall be forgiven him.''\endnote{James 5:14--15, Douay--Rheims, read 2026-07-26 in the tracked Challoner verse text. The Douay's ``priests'' renders \latin{presbyteros}, and Trent's third chapter takes the trouble to insist that the word means ordained ministers and not ``the elders by age, or the foremost in dignity amongst the people.'' The next verse (5:16, ``Confess therefore your sins one to another'') stands in the same paragraph and is not treated here.}

The first word of the prescription is \emph{sick}. Not dying; not \emph{in extremis}; sick. The apostolic text names an illness and promises two things---a raising up, and forgiveness if there is anything to forgive. The raising up is grammatically the principal promise and the forgiveness is conditional and secondary. Whatever the sacrament later came to be called, this is the text on which it rests, and it does not say what the later name says.

\subsection{What Trent taught, and what everyone thinks it taught}

The Council of Trent devoted the second half of its fourteenth session to the sacrament. Its four canons are definitive: that extreme unction is truly and properly a sacrament instituted by Christ and promulgated by James, not merely a rite received from the Fathers; that the anointing confers grace, remits sin, and comforts the sick, and has not ceased; that the Roman rite of it is not repugnant to James; that the ministers are ordained priests.\endnote{Council of Trent, session XIV (25 November 1551), canons on the sacrament of extreme unction, canons 1--4; Waterworth English and the Tauchnitz Latin, read 2026-07-26 in the full texts of those printings as digitized at the Internet Archive items registered in the source library for those editions. The canons carry anathemas and are definitive teaching. The doctrinal chapters that precede them are authoritative exposition of a lower register than the canons, though they are the same council's teaching and are cited by the Catechism at CCC 1523 (DS 1694, 1698).}

The doctrinal chapters state the effect: ``the thing here signified is the grace of the Holy Ghost, whose anointing cleanses away sins, if there be any still to be expiated, as also the remains of sins; and raises up and strengthens the soul of the sick person, by exciting in him a great confidence in the divine mercy; whereby the sick being supported, bears more easily the inconveniences and pains of his sickness, and more readily resists the temptations of the Devil\ldots{} and at times obtains bodily health, when expedient for the welfare of the soul.''\endnote{Trent, session XIV, doctrine on extreme unction, ch.~II, Waterworth English; Latin (\latin{cuius unctio delicta, si qua sint adhuc expianda, ac peccati reliquias abstergit, et aegroti animam alleviat et confirmat\ldots{} et sanitatem corporis interdum, ubi saluti animae expedierit, consequitur}) from the Tauchnitz printing. Read 2026-07-26 as above. The chapter's exposition follows the clauses of James 5:15 in order.} That is Trent's account of what the sacrament does, and every element of it is drawn from the verse.

Now the chapter that matters, and that almost nobody quotes. Chapter III, on the minister and on the time when the sacrament ought to be administered:

``It is also declared that this unction is to be applied \emph{to the sick}, but to those especially who lie in such danger as to seem to be about to depart this life: whence also it is called the sacrament of the departing. And if the sick should, after having received this unction, recover, they may again be aided by the succour of this sacrament, when they fall into another like danger of death.''\endnote{Trent, session XIV, doctrine on extreme unction, ch.~III, Waterworth English; the emphasis is this essay's. Latin from the Tauchnitz printing: ``\latin{Declaratur etiam, esse hanc unctionem infirmis adhibendam, illis vero praesertim, qui tam periculose decumbunt, ut in exitu vitae constituti videantur, unde et sacramentum exeuntium nuncupatur. Quod si infirmi post susceptam hanc unctionem convaluerint, iterum huius sacramenti subsidio iuvari poterunt, quum in aliud simile vitae discrimen inciderint.}'' Read 2026-07-26 as above. The construction is a preference within a class (\latin{infirmis\ldots{} illis vero praesertim}), not a restriction of the class.}

Read the Latin construction: the unction is to be given \latin{infirmis}---to the sick---\latin{illis vero praesertim}, but especially to those in danger. That is a preference within a class, not a restriction of it. Trent's own explanation of the name \emph{sacramentum exeuntium}, the sacrament of the departing, is offered as an account of where a customary title came from, immediately after stating the wider rule; and the very next sentence provides for the recipient's recovery and a second anointing later, which would be an odd provision in a sacrament reserved for the terminal.

So the received picture---that extreme unction was the last rite, restricted to the dying, until the Second Vatican Council relaxed it---does not survive a reading of the council that is supposed to have taught the restriction. Trent taught the wider rule and recorded the narrower custom. What happened afterwards was that the custom swallowed the rule, so thoroughly that a family calling for the priest often waited until the patient was unconscious, and the arrival of the priest became the household's signal that hope had been abandoned.

\subsection{The correction, and what it was correcting}

The Second Vatican Council's sentence is therefore a restoration and not an innovation, and it is worth noticing how carefully it is phrased: ``\,`Extreme unction,' which may also and more fittingly be called `anointing of the sick,' is not a sacrament for those only who are at the point of death. Hence, as soon as any one of the faithful begins to be in danger of death from sickness or old age, the fitting time for him to receive this sacrament has certainly already arrived.''\endnote{\emph{Sacrosanctum concilium} 73, Holy See English web text, read 2026-07-26. The constitution was promulgated 4 December 1963. Its next article, SC 74, prescribes that ``in addition to the separate rites for anointing of the sick and for viaticum, a continuous rite shall be prepared according to which the sick man is anointed after he has made his confession and before he receives viaticum''---which fixes the order of the three sacraments and is discussed below.} It does not say that the old name was wrong; it says the new one is more fitting. It does not announce a new discipline; it denies a misconception (``is not a sacrament for those only\ldots'') and then draws the practical consequence. The current law follows: anointing may be given to a member of the faithful ``who, having reached the use of reason, begins to be in danger due to sickness or old age,'' and may be repeated if the person recovers and again falls gravely ill, or if the same illness worsens.\endnote{\emph{Codex Iuris Canonici} (1983), can.~1004 §§1--2; Holy See English web text of the Code, read 2026-07-26. Compare CCC 1528: ``The proper time for receiving this holy anointing has certainly arrived when the believer begins to be in danger of death because of illness or old age,'' and CCC 1529 on repetition. The governing law is that of the Latin Church, in force since 27 November 1983; readers of the Eastern Catholic Churches are governed by the \emph{Codex Canonum Ecclesiarum Orientalium}, which this essay has not examined.}

The Catechism's account of the effects likewise stays close to Trent while pressing one point the council mentions only in passing: the sacrament gives ``strengthening, peace and courage to overcome the difficulties that go with the condition of serious illness or the frailty of old age\ldots{} against the temptations of the evil one, the temptation to discouragement and anguish in the face of death.''\endnote{CCC 1520, Holy See English web text, read 2026-07-26; the footnote is Hebrews 2:15, the verse on lifelong servitude through fear of death. CCC 1521--1523 add union with Christ's passion, an ecclesial grace, and the preparation for the final journey, the last quoting Trent's \latin{sacramentum exeuntium}.} That is the \emph{ars moriendi}'s five temptations, translated out of demonology into a sacramental grace, and attached to a rite rather than to a manual. The medieval tract handed a layman a script for the dying man's despair. The Church's answer to the same despair is a sacrament, and the sacrament came first.

\subsection{The last sacrament is not the one people think}

Here is the second place where the received picture is wrong, and this one is corrected in the plainest possible language.

``As the sacrament of Christ's Passover the Eucharist should always be the last sacrament of the earthly journey, the `viaticum' for `passing over' to eternal life.''\endnote{CCC 1517, Holy See English web text, read 2026-07-26; the paragraph also provides that the celebration ``can be preceded by the sacrament of Penance and followed by the sacrament of the Eucharist.'' CCC 1524--1525 develop the point: viaticum is ``the last sacrament of the Christian,'' and Penance, Anointing, and the Eucharist as viaticum together ``constitute at the end of Christian life `the sacraments that prepare for our heavenly homeland.'\,''} The order fixed by \emph{Sacrosanctum concilium} 74 is the same: confession, then anointing, then viaticum. And the law is emphatic about the last of the three: ``The Christian faithful who are in danger of death from any cause are to be nourished by holy communion in the form of Viaticum''; and, with a note of urgency rare in a code, ``Holy Viaticum for the sick is not to be delayed too long; those who have the care of souls are to be zealous and vigilant that the sick are nourished by Viaticum \emph{while fully conscious}.''\endnote{\emph{Codex Iuris Canonici} (1983), cann.~921 §1 and 922; Holy See English web text, read 2026-07-26; the emphasis is this essay's. Can.~921 §2 urges those in danger of death to receive again even if they have already communicated that day, and §3 recommends frequent communion, on separate days, while the danger lasts. The practice of reserving the Eucharist for the sick and carrying it to them is the subject of Trent, session XIII, ch.~VI and canon 7 (Waterworth English, read 2026-07-26), which anathematizes the denial that it is lawful ``that it be carried with honour to the sick.''}

\emph{While fully conscious.} The canon is legislating against a specific pastoral failure---the assumption that the sacraments are for the end, and that the end can be recognized in time. Everything in the Church's own arrangements pushes the opposite way: anoint early, communicate while he can still receive, and do not wait for the moment that the second section of this essay showed cannot be identified.

So the popular ordering is inverted at both ends. What most Catholics call ``the last rites'' is anointing, and anointing is not last and is not for the last. The last sacrament is the same one the Christian has been receiving all his life, given one more time under a name that means \emph{provision for a journey}.

\subsection{The Church's last words}

When there is nothing further to be done sacramentally, there is still something to be said, and the Church has words for it. The Catechism sets them down at the head of its treatment of life everlasting, after describing what has just been given: ``When the Church for the last time speaks Christ's words of pardon and absolution over the dying Christian, seals him for the last time with a strengthening anointing, and gives him Christ in viaticum as nourishment for the journey, she speaks with gentle assurance:''

``Go forth, Christian soul, from this world in the name of God the almighty Father, who created you, in the name of Jesus Christ, the Son of the living God, who suffered for you, in the name of the Holy Spirit, who was poured out upon you\ldots{} May you return to your Creator who formed you from the dust of the earth. May holy Mary, the angels, and all the saints come to meet you as you go forth from this life\ldots{} May you see your Redeemer face to face.''\endnote{CCC 1020, quoting the Prayer of Commendation from the Order of Christian Funerals; Holy See English web text, read 2026-07-26, and quoted in that official English rendering rather than in any translation composed for this essay. The ellipses mark omissions in the Catechism's own presentation and in this quotation; the full text appears in the ritual book. The prayer is the Latin rite's \latin{Proficiscere, anima christiana}, and the Catechism's placement of it---at the opening of the article on life everlasting, immediately before the particular judgment---is itself a piece of theology.}

Every claim this essay has examined is in that prayer, and none of it is argued. \emph{Go forth}: the separation is real, and it is a departure rather than an extinction. \emph{Return to your Creator who formed you from the dust of the earth}: Genesis 3:19, spoken not as a sentence but as an itinerary. \emph{Come to meet you}: the communion is not interrupted. \emph{See your Redeemer face to face}: \emph{Benedictus Deus}, in the vocative.

It is worth saying what kind of achievement that is. The prayer is addressed to a person who by definition cannot answer and may not be able to hear. It makes no promises about the individual's judgment---every clause is a petition or a commendation, not a verdict---and it does not console the family by pretending anything. What it does is name, out loud, in the room, what the tradition holds to be happening, at the exact moment when nobody can see it happening.

The verse this essay began with put the remembrance of the last things at the end of a list of things to do for other people: comfort the mourners, do not withhold kindness from the dead, be not slow to visit the sick. That order has turned out to be the right one. The doctrine of death is not, in the first place, a piece of information about the future; it is what makes it possible to stay in the room. A person who believes what has been set out here---that the dying man is a single nature coming apart, that the parting is a real deprivation and not a graduation, that what survives is not yet the man and will not be until the resurrection, that the accounting is immediate and the communion is not broken, and that nobody can locate the moment---has nothing to offer by way of explanation and every reason not to leave.

Of the four last things, this one can be attended. The remaining three are the subject of the volumes that follow: \emph{Judgment}, \emph{Hell}, and \emph{Heaven}. What death contributes to them is a definition, a boundary, and a warning about how much of what everyone believes on the subject was never taught.
